% sidenote_type=number（傍注，4.2.7）: pdfTeX検証
% 過去のpdfTeX不具合(a)の回帰フィクスチャ：number型\sidenoteが\marginpar
% 内部のボックス調整で「Dimension too large」（\jlreq@tempdima代入）．
% 真因は\set@fontsizeフックの\ptj@zwreg計算
% （\dimexpr\f@size\jlreq@jafontscale pt → "10"+"1"+"pt" と数字連結され
% 101pt化）で，\marginparの\@marginparresetが\normalsizeを呼んだ時点で
% 箱内の\zw実寸が101ptになり，続く\footnotesizeの
% \dimexpr115\jlreq@tempdima/170が溢れた．2段計算に修正済み（jlreq.cls d分岐）．
% （symbol型の\CJK@punctchar不具合(b)は sidenote-symbol.tex で検証．
% \jlreqsetupはプリアンブル専用なので別ファイルとした．）
% \footnote/\endnote/\marginparとの共存も併せて確認する．
\documentclass[pdflatex,sidenote_length=20zw]{jlreq}
\begin{document}
% --- sidenote_type=number（既定） ---
傍注（番号型）の呼び出し\sidenote{番号型の傍注本文。「括弧」や句読点，
English混植でもルールフックが効くことを見るための少し長い注。}の直後に本文が続く。

分離型：マークだけ出し\sidenotemark て注文は \sidenotetext{分離した番号型の傍注本文。}
あとから渡す形。オプション番号指定\sidenote[42]{明示番号の傍注。}も通す。

% --- 共存する脚注・後注・単体 marginpar ---
脚注\footnote{脚注本文。}と後注\endnote{後注本文。}，
単体の\marginpar{欄外の＊マークと\footnotesize 小さい文字}もサイズ追従する。

\theendnotes
\makeatletter
% 回帰ロック：マージン箱内でも\ptj@zwregが正しくサイズ追従すること
% （旧バグでは\marginpar箱内で101ptになっていた）
\setbox0=\vbox{\reset@font\footnotesize
  \ifdim\ptj@zwreg=8.00003pt\relax
    \typeout{SIDENOTE-JAREG-OK: \the\ptj@zwreg}%
  \else
    \errmessage{SIDENOTE-JAREG-NG: \the\ptj@zwreg\space (expected 8.00003pt)}%
  \fi
  X}%
\makeatother
\end{document}
